% ============================================================
%  BrushHeroes — Manual de Usuario  (DOCUMENTO DE MUESTRA)
%  Este archivo muestra el formato y estilo del manual final.
%  Las secciones completas se reemplazarán en el documento real.
% ============================================================
\documentclass[12pt, a4paper]{article}

% ── Paquetes ─────────────────────────────────────────────────
\usepackage[spanish]{babel}
\usepackage[utf8]{inputenc}
\usepackage[T1]{fontenc}
\usepackage{lmodern}
\usepackage{geometry}
\usepackage{graphicx}
\usepackage{float}
\usepackage{caption}
\usepackage{subcaption}
\usepackage{xcolor}
\usepackage{titlesec}
\usepackage{titling}
\usepackage{fancyhdr}
\usepackage{hyperref}
\usepackage{enumitem}
\usepackage{tcolorbox}
\usepackage{booktabs}
\usepackage{array}
\usepackage{parskip}

\geometry{
  top=2.5cm, bottom=2.5cm,
  left=3cm,  right=2.5cm
}

% ── Paleta de colores del juego ───────────────────────────────
\definecolor{heroBlue}{HTML}{2563EB}
\definecolor{heroGreen}{HTML}{22C55E}
\definecolor{heroOrange}{HTML}{F97316}
\definecolor{heroGray}{HTML}{64748B}
\definecolor{heroBg}{HTML}{F0F9FF}

% ── Tipografía de secciones ───────────────────────────────────
\titleformat{\section}
  {\Large\bfseries\color{heroBlue}}
  {\thesection.}{0.8em}{}[\titlerule]

\titleformat{\subsection}
  {\large\bfseries\color{heroBlue!80!black}}
  {\thesubsection.}{0.6em}{}

\titleformat{\subsubsection}
  {\normalsize\bfseries\color{heroGray}}
  {\thesubsubsection.}{0.5em}{}

% ── Cajas de aviso ────────────────────────────────────────────
\tcbuselibrary{skins, breakable}

\newtcolorbox{notebox}{
  colback=heroBg, colframe=heroBlue,
  fonttitle=\bfseries, title=Nota,
  left=6pt, right=6pt, top=4pt, bottom=4pt,
  arc=4pt
}

\newtcolorbox{tipbox}{
  colback=heroGreen!10, colframe=heroGreen!60!black,
  fonttitle=\bfseries, title=Consejo,
  left=6pt, right=6pt, top=4pt, bottom=4pt,
  arc=4pt
}

\newtcolorbox{warnbox}{
  colback=heroOrange!10, colframe=heroOrange,
  fonttitle=\bfseries, title=Importante,
  left=6pt, right=6pt, top=4pt, bottom=4pt,
  arc=4pt
}

% ── Encabezado y pie de página ────────────────────────────────
\pagestyle{fancy}
\fancyhf{}
\fancyhead[L]{\small\color{heroGray}BrushHeroes — Manual de Usuario}
\fancyhead[R]{\small\color{heroGray}\leftmark}
\fancyfoot[C]{\small\thepage}
\renewcommand{\headrulewidth}{0.4pt}

% ── Metadatos PDF ─────────────────────────────────────────────
\hypersetup{
  colorlinks=true,
  linkcolor=heroBlue,
  urlcolor=heroBlue,
  pdftitle={BrushHeroes - Manual de Usuario},
  pdfauthor={Equipo BrushHeroes}
}

% ─────────────────────────────────────────────────────────────
\begin{document}

% ── PORTADA ──────────────────────────────────────────────────
\begin{titlepage}
  \centering
  \vspace*{2cm}

  % Logo / imagen de portada
  \begin{figure}[H]
    \centering
    \includegraphics[width=0.45\textwidth]{img/portada_logo.png}
    % IMAGEN REQUERIDA: logo o splash art del juego con el nombre "BrushHeroes"
  \end{figure}

  \vspace{1cm}
  {\Huge\bfseries\color{heroBlue} BrushHeroes}\\[0.4cm]
  {\Large\color{heroGray} Manual de Usuario}\\[0.8cm]
  \rule{0.6\textwidth}{1pt}\\[0.8cm]
  {\large Videojuego educativo de higiene bucal para niños}\\[0.4cm]
  {\normalsize Realidad Aumentada · Android}\\[2cm]

  \vfill
  {\small Versión 1.0 \quad|\quad \today}\\[0.3cm]
  {\small Universidad \textbf{[nombre de tu universidad]}}\\[0.1cm]
  {\small Proyecto de Grado — Ingeniería \textbf{[tu carrera]}}
\end{titlepage}

% ── TABLA DE CONTENIDOS ───────────────────────────────────────
\tableofcontents
\newpage

% ─────────────────────────────────────────────────────────────
%  SECCIÓN 1 — INTRODUCCIÓN  (sección completa de muestra)
% ─────────────────────────────────────────────────────────────
\section{Introducción}

\textbf{BrushHeroes} es una aplicación móvil educativa diseñada para motivar a los niños a
mantener buenos hábitos de higiene bucal. A través de minijuegos interactivos y un sistema de
progreso gamificado, el jugador aprende la técnica correcta de cepillado de manera divertida.

\begin{figure}[H]
  \centering
  \includegraphics[width=0.55\textwidth]{img/intro_menu_general.png}
  % IMAGEN REQUERIDA: captura completa del menú principal (pestaña Minijuegos)
  \caption{Pantalla principal de BrushHeroes.}
  \label{fig:intro_menu}
\end{figure}

\subsection{¿A quién está dirigido?}

BrushHeroes está pensado para niños entre 4 y 10 años, aunque puede ser utilizado por
cualquier persona. Se recomienda la supervisión de un adulto para la configuración inicial.

\subsection{Características principales}

\begin{itemize}[leftmargin=1.5cm]
  \item \textbf{Minijuego de Realidad Aumentada (AR):} usa la cámara frontal del dispositivo
        para superponer bacterias y un cepillo virtual sobre la cara del jugador, enseñando
        las zonas de cepillado.
  \item \textbf{Minijuego de Cepillado:} práctica interactiva de la técnica de cepillado.
  \item \textbf{Minijuego de Hilo Dental:} aprende a usar el hilo dental correctamente.
  \item \textbf{Sistema de racha:} registra cuántos días consecutivos el jugador ha jugado.
  \item \textbf{Estrellas y misiones:} recompensas por completar sesiones matutinas y vespertinas.
  \item \textbf{Calendario de actividad:} visualiza el historial de juego del mes.
\end{itemize}

\begin{notebox}
  Esta aplicación \textbf{no reemplaza} la instrucción de un profesional de salud bucal.
  Su propósito es reforzar hábitos de forma lúdica.
\end{notebox}

% ─────────────────────────────────────────────────────────────
%  SECCIÓN 2 — REQUISITOS  (sección completa de muestra)
% ─────────────────────────────────────────────────────────────
\section{Requisitos del sistema}

\begin{table}[H]
  \centering
  \caption{Requisitos mínimos del dispositivo.}
  \label{tab:requisitos}
  \begin{tabular}{>{\bfseries}p{4cm} p{8cm}}
    \toprule
    Componente & Requisito \\
    \midrule
    Sistema operativo  & Android 9.0 (API 28) o superior \\
    Procesador         & Octa-core 2.0\,GHz o superior \\
    Memoria RAM        & 3\,GB mínimo (4\,GB recomendado) \\
    Almacenamiento     & 300\,MB de espacio libre \\
    Cámara frontal     & Requerida (resolución mínima 5\,MP) \\
    Soporte AR         & ARCore compatible (ver lista oficial de Google) \\
    Conexión a red     & No requerida para jugar \\
    \bottomrule
  \end{tabular}
\end{table}

\begin{warnbox}
  El \textbf{minijuego de Realidad Aumentada} requiere que el dispositivo sea compatible con
  \textbf{Google ARCore}. Puede consultar la lista de dispositivos compatibles en
  \href{https://developers.google.com/ar/devices}{developers.google.com/ar/devices}.
\end{warnbox}

% ─────────────────────────────────────────────────────────────
%  SECCIÓN 3 — INSTALACIÓN  (esqueleto — solo muestra estructura)
% ─────────────────────────────────────────────────────────────
\section{Instalación}

\subsection{Instalación desde archivo APK}

\begin{enumerate}[leftmargin=1.5cm]
  \item Descargue el archivo \texttt{BrushHeroes.apk} desde el enlace proporcionado por su
        institución educativa.
  \item En su dispositivo Android, vaya a \textbf{Ajustes → Seguridad} y active la opción
        \textit{«Instalar apps de fuentes desconocidas»}.
  \item Abra el archivo \texttt{.apk} descargado y siga las instrucciones en pantalla.
  \item Una vez instalada, la aplicación aparecerá en su menú de aplicaciones con el ícono
        de BrushHeroes.
\end{enumerate}

\begin{figure}[H]
  \centering
  \begin{subfigure}[b]{0.3\textwidth}
    \includegraphics[width=\textwidth]{img/install_step1.png}
    % IMAGEN REQUERIDA: ajustes de Android — activar fuentes desconocidas
    \caption{Activar fuentes desconocidas.}
  \end{subfigure}
  \hfill
  \begin{subfigure}[b]{0.3\textwidth}
    \includegraphics[width=\textwidth]{img/install_step2.png}
    % IMAGEN REQUERIDA: diálogo de instalación del APK
    \caption{Instalar el APK.}
  \end{subfigure}
  \hfill
  \begin{subfigure}[b]{0.3\textwidth}
    \includegraphics[width=\textwidth]{img/install_step3.png}
    % IMAGEN REQUERIDA: ícono de la app en el lanzador de Android
    \caption{App instalada.}
  \end{subfigure}
  \caption{Proceso de instalación paso a paso.}
  \label{fig:instalacion}
\end{figure}

% ─────────────────────────────────────────────────────────────
%  SECCIÓN 4 — MENÚ PRINCIPAL  (esqueleto)
% ─────────────────────────────────────────────────────────────
\section{Menú principal}

\textit{[Contenido completo en el documento final]}

\begin{figure}[H]
  \centering
  \includegraphics[width=0.5\textwidth]{img/menu_minijuegos_page.png}
  % IMAGEN REQUERIDA: captura de la pestaña "Minijuegos" del menú
  \caption{Pestaña Minijuegos — vista completa.}
  \label{fig:menu_minijuegos}
\end{figure}

% ─────────────────────────────────────────────────────────────
%  SECCIÓN 5 — MINIJUEGO AR  (subsección de muestra)
% ─────────────────────────────────────────────────────────────
\section{Minijuego AR: Cepillado con Realidad Aumentada}

\subsection{Acceso al minijuego}

Desde la pantalla principal, toque la tarjeta \textbf{«Cepillado AR»} y luego el botón
\textbf{«Jugar»}. La aplicación abrirá la cámara frontal y comenzará a buscar su rostro.

\subsection{Detección de rostro}

\begin{figure}[H]
  \centering
  \begin{subfigure}[b]{0.44\textwidth}
    \includegraphics[width=\textwidth]{img/ar_buscando_rostro.png}
    % IMAGEN REQUERIDA: pantalla de inicio del AR con overlay "Buscando rostro..."
    \caption{Buscando rostro\ldots}
  \end{subfigure}
  \hfill
  \begin{subfigure}[b]{0.44\textwidth}
    \includegraphics[width=\textwidth]{img/ar_rostro_detectado.png}
    % IMAGEN REQUERIDA: pantalla de inicio con cara detectada y botones activos
    \caption{Rostro detectado — botones habilitados.}
  \end{subfigure}
  \caption{Estados de la pantalla de inicio del minijuego AR.}
  \label{fig:ar_face_detection}
\end{figure}

Apunte la cámara frontal hacia su rostro. Cuando la aplicación detecte su cara, los botones
de selección de modo se habilitarán automáticamente.

\begin{tipbox}
  Para una mejor detección, asegúrese de estar en un lugar bien iluminado y de tener el
  rostro completamente visible en la pantalla.
\end{tipbox}

\subsection{Selección de modo de juego}

\begin{figure}[H]
  \centering
  \includegraphics[width=0.5\textwidth]{img/ar_seleccion_modo.png}
  % IMAGEN REQUERIDA: pantalla de inicio con los dos botones: "Dinámico" y "Solo guía"
  \caption{Pantalla de selección de modo.}
  \label{fig:ar_modo}
\end{figure}

BrushHeroes AR ofrece dos modos de juego:

\begin{itemize}[leftmargin=1.5cm]
  \item \textbf{Modo Dinámico:} el jugador arrastra su dedo por la pantalla para mover el
        cepillo virtual sobre las bacterias. Cuenta con un cronómetro de 3 minutos.
  \item \textbf{Modo Solo Guía:} el cepillo se mueve automáticamente, mostrando la técnica
        correcta. Ideal para niños más pequeños o para aprender la rutina. Sin límite de tiempo.
\end{itemize}

\subsubsection{Modo Dinámico — ¿Cómo jugar?}

\begin{figure}[H]
  \centering
  \includegraphics[width=0.55\textwidth]{img/ar_dinamico_gameplay.png}
  % IMAGEN REQUERIDA: gameplay modo dinámico — cara con bacterias visibles,
  %   cepillo virtual posicionado, HUD en la parte superior
  \caption{Modo Dinámico en acción.}
  \label{fig:ar_dinamico}
\end{figure}

\begin{enumerate}[leftmargin=1.5cm]
  \item Las bacterias aparecerán sobre distintas zonas de su boca.
  \item \textbf{Arrastre su dedo} por la pantalla para mover el cepillo virtual.
  \item Pase el cepillo sobre todas las bacterias de la zona activa para limpiarla.
  \item Al limpiar una zona, aparecerá una notificación y pasará a la siguiente.
  \item Complete las 6 zonas antes de que se acabe el tiempo para ganar.
\end{enumerate}

% ─────────────────────────────────────────────────────────────
%  SECCIÓN 6 — PROGRESO (esqueleto)
% ─────────────────────────────────────────────────────────────
\section{Sistema de progreso}

\textit{[Contenido completo en el documento final]}

% ─────────────────────────────────────────────────────────────
%  SECCIÓN 7 — PREGUNTAS FRECUENTES  (muestra parcial)
% ─────────────────────────────────────────────────────────────
\section{Preguntas frecuentes}

\subsection*{¿Por qué la cámara no detecta mi rostro?}
Asegúrese de estar en un lugar bien iluminado. Si el problema persiste, verifique que su
dispositivo sea compatible con Google ARCore.

\subsection*{¿Se pierde mi progreso si desinstalo la app?}
Sí. El progreso se guarda localmente en el dispositivo. Desinstalar la aplicación eliminará
todos los datos guardados.

\subsection*{¿Necesito conexión a internet?}
No. BrushHeroes funciona completamente sin conexión una vez instalado.

\subsection*{¿Puedo usar auriculares?}
Sí. El juego tiene efectos de sonido y música de fondo. El uso de auriculares es opcional.

\vfill
\begin{center}
  \small\color{heroGray}
  BrushHeroes · Proyecto de Grado · \the\year \\
  Documento generado para revisión de formato
\end{center}

\end{document}
